\section{Introduction}
Lattice QCD simulations predict an ever increasing number of observables
that are relevant to particle and hadron physics phenomenology.
These results are usually extracted from expectation values
of $n$-point functions at large Euclidean time separations.
Due to the decay of these functions with time, statistical noise
over signal ratios increase exponentially as time separations
are taken large (with the notable exception of pseudoscalar
mesons at zero momentum).
Fortunately, there exists some freedom in the construction of interpolators
for the creation of mesonic and baryonic states. Employing interpolators
that resemble the spatial structure of
the ground state wave function enables asymptotic results to be extracted at
time separations where the signal over noise ratio is still large.

Many applications nowadays demand hadrons that carry momentum.
For instance, pushing the calculation of semileptonic decay form
factors for $B\rightarrow\pi\ell\bar{\nu}_{\ell}$ or
$\Lambda_b\rightarrow p\ell\bar{\nu}_{\ell}$~\cite{Detmold:2015aaa}
towards small virtualities requires spatial momenta of the
size of the mass difference between the $B$ meson and
the pion or between the $\Lambda_b$ baryon and the proton, respectively.
Another important type of application are parton distribution functions (PDFs) and their generalizations, in particular transverse momentum dependent parton distribution functions (TMDs), or Wigner distributions as a whole. 
Also these quantities are extracted from matrix elements of the type $ \bra{H(\vect{p}')} \oper \ket{H(\vect{p})}$, where $H(\vect{p})$
is a hadron state with momentum $\vect{p}$.
The operator $\oper$ cannot have an extent in Minkowski-time on the Euclidean
lattice. Therefore, in order to extrapolate to light front kinematics for an inherently non-local operator,
one has to work in a frame of reference where hadrons carry high spatial momenta $\vect{p}$, $\vect{p}'$.
This fact has been known for quite some time in the context of lattice calculations of TMDs \cite{Hagler:2009mb,Musch:2010ka,Musch:2011er} 
and is illustrated very clearly in recent work on these distributions in the pion~\cite{Engelhardt:2015xja}.
For the same reason fast hadrons on the lattice are highly desirable in a new scheme proposed 
to relate quasi parton distributions to light front
distributions in a controlled manner \cite{Ji:2013dva,Ji:2015jwa}.
First lattice
computations in this direction have
started~\cite{Lin:2014zya,Alexandrou:2015rja}.
Earlier suggestions to compute quasi distribution amplitudes
in position space~\cite{Aglietti:1998ur,Abada:2001if,Braun:2007wv}
equally require
pions or nucleons at high momenta. Unfortunately,
up to now no satisfactory techniques for hadrons carrying high momenta
existed to suppress excited
state contributions.

To be more specific, a two-point function
is given as
\begin{equation}
\label{eq:twopoint}
C_H(t)=\langle H(t)H^{\dagger}(0)\rangle
=\sum_{j>0}\frac{|\langle j|\hat{H}^{\dagger}|0\rangle|^2}{2E_{H,j}}e^{-E_{H,j}t}\,,
\end{equation}
where $E_{H,j}$ denotes the $j$th energy level within the tower
of states created by the interpolator
$H^{\dagger}$. Obviously, the contribution of the $j$th excited
state is suppressed relative to the ground state
not only by $\exp[-(E_{H,j}-E_{H,1})t]$ but also
by the ratio
$|\langle j|\hat{H}^{\dagger}|0\rangle|^2/|\langle 1|\hat{H}^{\dagger}|0\rangle|^2$:
increasing the ground state overlap factor
$|\langle 1|\hat{H}^{\dagger}|0\rangle|$,
relative to $|\langle j|\hat{H}^{\dagger}|0\rangle|$ for $j>1$,
results in an additional suppression of excitations.

Reducing excited state overlaps by employing extended
interpolators was first pursued in computations of the
glueball spectrum. In this case the gauge links within the corresponding
interpolators can be iteratively ``APE smeared''~\cite{Falcioni:1984ei},
``fuzzed''~\cite{Teper:1987wt} or ``HYP smeared''~\cite{Hasenfratz:2001hp},
to better approximate the (smooth) ground state wave function, see
also Ref.~\cite{Gupta:1990ek}.
This gauge link smearing was subsequently generalized
to iterative smearing of quark fields within interpolators
that create mesonic and baryonic states, in particular gauge covariant Wuppertal
(i.e.~Gau\ss{}) smearing~\cite{Gusken:1989ad,Gusken:1989qx,Alexandrou:1990dq},
hydrogen-like smearing~\cite{Alexandrou:1990dq}
as well as Jacobi smearing~\cite{Collins:1992rp,Allton:1993wc}.
Additionally, in Refs.~\cite{Gusken:1989qx,Alexandrou:1990dq,Bali:2005fu}
APE smearing was employed for the spatial gauge transporters within
the quark smearing while in Ref.~\cite{Bali:2005fu}
linear combinations of different levels of Wuppertal smearing
were utilized.

Since Gaussian smearing functions may not be optimal for creating, e.g.,
$p$-waves, even when adding derivatives to the interpolator,
iterative smearing was later-on combined with
displaced quark sources (fuzzing) in Ref.~\cite{Lacock:1994qx},
a generalization of which was suggested in Ref.~\cite{Boyle:1999gx}.
Finally, in Ref.~\cite{vonHippel:2013yfa} ``free form smearing'',
folding Gaussian smearing with an arbitrary function in a gauge
covariant way, was invented. Preceding and in parallel to
gauge covariant iterative smearing functions, gauge fixed sources have
been utilized: wall sources for zero~\cite{Bacilieri:1988fh} and
non-zero momentum~\cite{Detmold:2012wc}, box~\cite{Allton:1990qg}
sources, Gaussian ``shell sources''~\cite{DeGrand:1990dz}
and sources with nodes~\cite{DeGrand:1991ng}.
These gauge fixed methods and free form smearing share the disadvantage
that smearing the sink requires all quark positions to
be summed over individually, turning this prohibitively expensive.
Having identical source and sink interpolators,
however, is very desirable as only this guarantees the
positivity of the coefficients of the spectral decomposition
Eq.~\eqref{eq:twopoint} and thus the convexity of two-point functions.
For completeness, we also mention the ``distillation'' (or
Laplacian-Heaviside) method of Ref.~\cite{Peardon:2009gh}
since this is closely related to gauge covariant smearing.

Large momenta increase the energy of the state and result in faster
decaying two- and three-point functions and, therefore, in inferior noise to
signal ratios. Moreover, as we shall see,
excited state suppression becomes far less effective when
using conventional quark smearing methods.
Some attempts have been
made~\cite{Roberts:2012tp,DellaMorte:2012xc}
to introduce an anisotropy into Wuppertal
smearing~\cite{Gusken:1989ad,Gusken:1989qx}, aiming at
Lorentz contracting the interpolating wave function
according to the boost factor $1/\gamma=m/E(\vect{p})$,
along the direction of the spatial
momentum $\vect{p}$. However,
this did not result in the ground state enhancement that one would have
hoped for. Here we will argue and demonstrate that to achieve satisfactory
results at high momenta, additional phase factors need to be incorporated
into quark smearing functions.

This article is organized as follows.
First, in Sec.~\ref{sec:basic}, we discuss the basic idea behind the new
class of smearing functions that we introduce.
Then, in Sec.~\ref{sec:wuppertal} we are more specific, modifying
Wuppertal smearing as a generic example and suggest further improvements.
In Sec.~\ref{sec:simulate} we discuss our simulation
parameters and expectations for
the nucleon and pion energies. After the stage is set,
in Sec.~\ref{sec:results} we investigate the feasibility of the method
in a realistic numerical study, optimize the smearing parameters and pursue a
comparison between the new and the conventional method. Finally,
we study the pion and nucleon dispersion relations, before we conclude.

